%==============================================================================
%  IMPLEMENTACIÓN EN PYTORCH — PARA DUMMIES
%  Cómo está implementado en el repositorio el MPNN, la entalpía y sus tablas,
%  el GENERIC dentro de la red, el entrenamiento push-forward (K) y la pérdida.
%
%  Documento hermano de "generic_entalpico_para_dummies": aquel explica QUÉ
%  queremos hacer; éste explica CÓMO está implementado, ecuación <-> código.
%
%  Nivel de partida del lector: sabe entrenar una red en PyTorch (init, forward,
%  loss, bucle de entrenamiento y validación; nivel "MNIST").
%==============================================================================
\documentclass[11pt]{article}

\usepackage[utf8]{inputenc}
\usepackage[T1]{fontenc}
\usepackage[spanish,es-noshorthands]{babel}
\usepackage[margin=1in]{geometry}
\usepackage{amsmath,amssymb,amsthm,mathtools}
\usepackage{bm}
\usepackage{graphicx}
\usepackage{booktabs}
\usepackage{array}
\usepackage{xcolor}
\usepackage{enumitem}
\usepackage{listings}
\usepackage[hidelinks,unicode]{hyperref}

% --- macros matemáticas -------------------------------------------------------
\newcommand{\dd}{\mathrm{d}}
\newcommand{\Lw}{\bm{L}_{w}}
\newcommand{\cpapp}{c_p^{\mathrm{app}}}
\newcommand{\cpsens}{c_p^{\mathrm{sens}}}
\newcommand{\Tinf}{T_{\infty}}

% --- cajas pedagógicas --------------------------------------------------------
\theoremstyle{definition}
\newtheorem{idea}{Idea clave}[section]
\theoremstyle{plain}

\newcommand{\dummies}[1]{%
  \par\smallskip\noindent
  \colorbox{yellow!18}{\parbox{\dimexpr\linewidth-2\fboxsep}{%
  \small\textbf{En cristiano:} #1}}\par\smallskip}

% --- estilo de código ---------------------------------------------------------
\definecolor{codebg}{rgb}{0.96,0.96,0.99}
\definecolor{codekw}{rgb}{0.10,0.10,0.60}
\definecolor{codecom}{rgb}{0.28,0.45,0.30}
\definecolor{codestr}{rgb}{0.55,0.15,0.15}
\definecolor{srcblue}{rgb}{0.12,0.25,0.55}

\lstdefinestyle{py}{
  language=Python,
  basicstyle=\ttfamily\footnotesize,
  keywordstyle=\color{codekw}\bfseries,
  commentstyle=\color{codecom}\itshape,
  stringstyle=\color{codestr},
  backgroundcolor=\color{codebg},
  frame=single, framerule=0.4pt, rulecolor=\color{srcblue!40},
  numbers=left, numberstyle=\tiny\color{gray}, numbersep=7pt,
  showstringspaces=false, breaklines=true, breakatwhitespace=true,
  tabsize=4, columns=fullflexible, keepspaces=true,
  literate=%
    {μ}{{$\mu$}}1 {Δ}{{$\Delta$}}1 {ρ}{{$\rho$}}1 {σ}{{$\sigma$}}1
    {≥}{{$\geq$}}1 {≤}{{$\leq$}}1 {→}{{$\rightarrow$}}1 {⇒}{{$\Rightarrow$}}1
    {·}{{$\cdot$}}1 {α}{{$\alpha$}}1 {∇}{{$\nabla$}}1 {≈}{{$\approx$}}1
    {Σ}{{$\Sigma$}}1 {τ}{{$\tau$}}1 {ε}{{$\varepsilon$}}1 {²}{{\textsuperscript{2}}}1
    {⟺}{{$\Leftrightarrow$}}1 {é}{{\'e}}1 {á}{{\'a}}1 {í}{{\'i}}1 {ó}{{\'o}}1
    {ú}{{\'u}}1 {ñ}{{\~n}}1
}
\lstset{style=py}

% cabecera "archivo — descripción" antes de cada caja de código
\newcommand{\srcref}[1]{\par\smallskip\noindent
  {\footnotesize\ttfamily\textcolor{srcblue}{$\blacksquare$\ #1}}\par\vspace{1pt}}

\title{\bfseries Implementación en PyTorch \emph{para dummies}\\[4pt]
  \large Del MPNN, la entalpía y el GENERIC a las líneas de código\\
  que lo hacen: una guía del repositorio, ecuación $\leftrightarrow$ archivo}
\author{}
\date{}

\begin{document}
\maketitle

\begin{abstract}
\noindent
Este es el documento \textbf{de implementación}. El documento hermano
(\emph{GENERIC entálpico para dummies}) explicaba \emph{qué} queremos hacer y
\emph{por qué}; aquí explicamos \emph{cómo está construido} en el repositorio,
relacionando cada pieza de la formulación con \textbf{las líneas concretas de
PyTorch que la implementan}, indicando siempre el \textbf{archivo y las líneas}.
Suponemos que sabes entrenar una red en PyTorch (un MNIST: \texttt{\_\_init\_\_},
\texttt{forward}, \emph{loss}, bucle de entrenamiento y validación). A partir de
ahí recorremos: el mapa de la librería y el flujo de datos; el contrato de grafos
(las 12 \emph{features}); la arquitectura \emph{Encoder--Processor--Decoder}; la
normalización; el \emph{head} GENERIC entálpico con sus conductancias aprendidas y
las \textbf{tablas que pasan de entalpía a temperatura}; el entrenamiento
\emph{push-forward} con $K$ pasos y su pérdida; y la validación por \emph{rollout}.
Al final tendrás un mapa mental completo del repo y podrás responder casi cualquier
pregunta sobre él. Las cajas amarillas resumen ``en cristiano''; las cajas azules
son código \textbf{real} del repositorio (a veces abreviado, marcado con
\texttt{...}).
\end{abstract}

\tableofcontents
\bigskip

%==============================================================================
\section{El mapa de la librería y el flujo de datos}
\label{sec:mapa}
%==============================================================================

Antes de mirar una sola línea, conviene saber \textbf{qué hace cada carpeta} y
\textbf{cómo fluye un dato} desde la física hasta una figura. Todo el código vive
en \texttt{src/}:

\begin{center}\small
\begin{tabular}{@{}p{0.34\linewidth}p{0.60\linewidth}@{}}
\toprule
Archivo & Responsabilidad \\
\midrule
\texttt{src/simulation/thermal\_solver.py} & El \textbf{solver FEM}: resuelve la
ecuación del calor (Goldak, $\cpapp$, $\theta$-método + Picard) y define
\texttt{MaterialProperties} y \texttt{SimulationResult}. Es la ``verdad''. \\
\texttt{src/simulation/generate\_dataset.py} & Genera el corpus: muestrea
parámetros y lanza el solver en paralelo $\rightarrow$ \texttt{.npz}. \\
\texttt{src/data/graph\_builder.py} & Convierte cada \texttt{.npz} en
\textbf{grafos PyG} (12 features de nodo, 3 de arista, target $\Delta T$);
normalización; el \texttt{Dataset} de PyTorch. \\
\texttt{src/models/meshgraphnet.py} & La \textbf{red}: Encoder/Processor/Decoder y
los \emph{heads} GENERIC (incluido el \textbf{entálpico} con su tabla). \\
\texttt{src/training/utils.py} & \emph{Split} por simulación, ruido de
entrenamiento, y las \textbf{ventanas} \texttt{WindowedSubset}/\texttt{collate\_windows}
del push-forward. \\
\texttt{src/training/train.py} & El \textbf{bucle de entrenamiento}: push-forward
$K$, pérdida, optimizador, \emph{scheduler}, \emph{checkpoints}, validación. \\
\texttt{src/training/rollout.py} & El \textbf{rollout autoregresivo} (inferencia)
y su exportación; es la métrica de validación. \\
\texttt{src/training/spiral\_*.py}, \texttt{val\_analysis.py} & Evaluación
fuera-de-distribución (la espiral) y análisis/figuras. \\
\bottomrule
\end{tabular}
\end{center}

El \textbf{flujo de datos} de punta a punta:
\begin{center}\small\ttfamily
parámetros $\rightarrow$ [thermal\_solver] $\rightarrow$ sim\_*.npz (campos T en el tiempo)\\
$\rightarrow$ [graph\_builder] $\rightarrow$ data\_\{sim\}\_\{t\}.pt (grafos) + stats.pt\\
$\rightarrow$ [train.py] $\rightarrow$ best\_model.pt / last\_model.pt\\
$\rightarrow$ [rollout.py / spiral\_*] $\rightarrow$ .npz + .xdmf + figuras
\end{center}

\dummies{Cinco etapas: el solver hace la física (\texttt{.npz}); el
\emph{graph\_builder} la vuelve grafos (\texttt{.pt}); \emph{train.py} aprende y
guarda pesos; \emph{rollout.py} despliega la red sobre una simulación entera y la
compara con la verdad. Cada flecha es un archivo. Quédate con esto y ya sabes
``dónde mirar'' para cualquier pregunta.}

\begin{idea}[El contrato que une todo]
Tres constantes amarran toda la librería y aparecen una y otra vez:
\textbf{12 features de nodo} (orden fijo en \texttt{NODE\_FEATURE\_NAMES}),
\textbf{3 de arista} \texttt{[dx,dy,dist]}, y el \textbf{target} $\bm
y=\Delta T=T^{t+1}-T^{t}$. El modelo \emph{solo} predice $\Delta T$; todo lo demás
(features, normalización, rollout) gira alrededor de ese contrato.
\end{idea}

%==============================================================================
\section{El contrato de datos: qué es ``un grafo'' aquí}
\label{sec:datos}
%==============================================================================

\subsection{Las 12 features de nodo y el target}

Cada par de fotogramas consecutivos $(T^t,T^{t+1})$ de una simulación se convierte
en \textbf{un grafo}: los nodos de la malla son nodos del grafo, las aristas de la
malla son aristas (en ambos sentidos), las features describen el estado en $t$ y el
target es el incremento de temperatura. El orden de las 12 columnas es \textbf{fijo}
y todo el código lo referencia por nombre/índice:

\srcref{src/data/graph\_builder.py — l.48--73 (layout de features)}
\begin{lstlisting}
NODE_FEATURE_NAMES = [
    "T",            # 0  temperatura actual  <- la UNICA que se "rolla"
    "q_goldak",     # 1  fuente Goldak analitica en el nodo
    "dx_local",     # 2  coord. relativa co-movil (tangente)
    "dy_local",     # 3  coord. relativa co-movil (normal)
    "net_power",    # 4  parametro de proceso: eta * P
    "speed",        # 5  velocidad de soldadura v
    "bc_interior",  # 6  one-hot tipo de nodo
    "bc_dirichlet", # 7  one-hot tipo de nodo
    "bc_robin",     # 8  one-hot tipo de nodo
    "h_conv",       # 9  coef. de conveccion local
    "emissivity",   # 10 emisividad local
    "T_inf",        # 11 temperatura ambiente local
]
NORMALIZE_MASK = np.ones(NUM_NODE_FEATURES, dtype=bool)
NORMALIZE_MASK[[6, 7, 8]] = False   # los one-hot NO se normalizan
\end{lstlisting}

Dos decisiones de diseño con consecuencias enormes:
\begin{itemize}[leftmargin=1.4em]
  \item \textbf{Solo features relativas}: no hay coordenadas absolutas. La posición
  del arco entra como \texttt{q\_goldak} (la fuente evaluada en el nodo) y como las
  coordenadas co-móviles \texttt{[dx\_local, dy\_local]} en el marco
  (tangente, normal) de la antorcha. La geometría de la malla entra solo por las
  aristas. Así el modelo no memoriza posiciones: aprende física local.
  \item \textbf{La columna 0 (T) es la única que se realimenta} en el rollout (lo
  veremos en §\ref{sec:rollout}); las otras 11 se conocen de antemano para todos
  los pasos.
\end{itemize}

El target se construye así (resta de fotogramas), y cada grafo se empaqueta como un
\texttt{Data} de PyTorch Geometric:

\srcref{src/data/graph\_builder.py — build\_graph\_sequence, l.246--260 (abreviado)}
\begin{lstlisting}
for t in range(n_pairs):
    x_np = build_node_features(result, t, speeds[t], node_type, bc_values, goldak)
    y_np = (result.temperature[t+1] - result.temperature[t]).reshape(-1, 1)  # dT
    data = Data(
        x=torch.as_tensor(x_np, dtype=torch.float32),         # (N, 12)
        edge_index=edge_index,                                 # (2, 2E) bidireccional
        edge_attr=edge_attr,                                   # (2E, 3) [dx,dy,dist]
        y=torch.as_tensor(y_np, dtype=torch.float32),          # (N, 1) = dT
        pos=pos,                                               # (N,2) SOLO para dibujar
    )
\end{lstlisting}

\dummies{Un grafo $=$ una foto del estado en $t$ (12 números por nodo) $+$ la
respuesta correcta $\Delta T$ por nodo. ``Entrenar'' es enseñar a la red a predecir
ese $\Delta T$. Fíjate que \texttt{pos} (coordenadas) se guarda aparte y nunca entra
en \texttt{x}: es solo para ParaView.}

\subsection{Las aristas y la fuente Goldak}

Las aristas se construyen una vez por malla (las tres aristas de cada triángulo,
sin duplicados, en ambos sentidos), y sus features son el desplazamiento
$[\Delta x,\Delta y,\lVert\bm u_{ij}\rVert]$. La \texttt{q\_goldak} de cada nodo la
calcula la \emph{misma} función del solver (\texttt{goldak\_flux}), multiplicada por
un ``power\_ratio'' que vale 1 soldando y 0 en el enfriamiento (cuando el arco se
apaga). No nos detenemos: lo importante es que \textbf{todas las features
dependientes del tiempo se calculan con el mismo código en entrenamiento y en
inferencia} (paridad), lo que evita el clásico \emph{train/test skew}.

%==============================================================================
\section{La arquitectura MPNN en PyTorch}
\label{sec:mpnn}
%==============================================================================

El modelo es el \emph{MeshGraphNet} (Pfaff et al.) clásico:
\textbf{Encoder $\rightarrow$ Processor $\rightarrow$ Decoder}. Todo está en
\texttt{src/models/meshgraphnet.py}.

\subsection{El ladrillo: un MLP configurable}

Todas las subredes son MLPs construidos por la misma fábrica: \texttt{num\_mlp\_layers}
capas ocultas de ancho \texttt{hidden\_dim} con activación, y al final una proyección
lineal (más un \texttt{LayerNorm} opcional, estándar en MeshGraphNets para estabilizar
las latentes; el decoder es la excepción, sin LayerNorm, para poder emitir valores
físicos sin acotar).

\srcref{src/models/meshgraphnet.py — build\_mlp, l.140--149}
\begin{lstlisting}
layers = []
dim = in_dim
for _ in range(num_hidden_layers):
    layers.append(nn.Linear(dim, hidden_dim)); layers.append(activation())
    dim = hidden_dim
layers.append(nn.Linear(dim, out_dim))
if layer_norm:
    layers.append(nn.LayerNorm(out_dim))
return nn.Sequential(*layers)
\end{lstlisting}

\subsection{Encoder: subir features a la latente}

Dos MLPs independientes llevan las features crudas de nodo (12-d) y de arista (3-d)
a un espacio latente compartido de ancho $d=128$. Esto es la ecuación
$\bm v_i^{(0)}=\mathrm{MLP}_{\mathrm{enc}}^{v}(\bm x_i)$,
$\bm\epsilon_{ij}^{(0)}=\mathrm{MLP}_{\mathrm{enc}}^{e}(\bm e_{ij})$.

\srcref{src/models/meshgraphnet.py — Encoder.forward, l.170--173}
\begin{lstlisting}
def forward(self, x, edge_attr):
    return self.node_mlp(x), self.edge_mlp(edge_attr)   # (N,d), (2E,d)
\end{lstlisting}

\subsection{Processor: el paso de mensajes (el corazón)}

El Processor es una \textbf{pila de $L=8$ bloques} \texttt{GraphNetBlock}. Cada
bloque hace tres cosas, con \textbf{conexiones residuales} (sumar la entrada) para
que el gradiente sobreviva la pila profunda y el BPTT del push-forward:

\begin{enumerate}[leftmargin=1.4em]
  \item \textbf{Actualiza cada arista} a partir de sí misma y de sus dos extremos:
  $\bm\epsilon_{ij}\leftarrow\bm\epsilon_{ij}+\mathrm{MLP}_e([\bm\epsilon_{ij},\bm v_i,\bm v_j])$.
  \item \textbf{Agrega} los mensajes en el nodo receptor (suma): aquí es donde el
  ``message passing'' mueve información entre vecinos. Se hace con
  \texttt{scatter} (una suma segmentada vectorizada).
  \item \textbf{Actualiza cada nodo} con su estado y la suma de mensajes:
  $\bm v_j\leftarrow\bm v_j+\mathrm{MLP}_v([\bm v_j,\bm m_j])$.
\end{enumerate}

\srcref{src/models/meshgraphnet.py — GraphNetBlock.forward, l.209--224}
\begin{lstlisting}
src, dst = edge_index[0], edge_index[1]          # emisor i, receptor j
# 1) update de arista (residual)
edge_inputs = torch.cat([edge_attr, x[src], x[dst]], dim=-1)
edge_attr = edge_attr + self.edge_mlp(edge_inputs)
# 2) agregacion en el receptor (suma segmentada vectorizada)
aggregated = scatter(edge_attr, dst, dim=0, dim_size=x.size(0), reduce="sum")
# 3) update de nodo (residual)
node_inputs = torch.cat([x, aggregated], dim=-1)
x = x + self.node_mlp(node_inputs)
return x, edge_attr
\end{lstlisting}

\dummies{\texttt{scatter(..., dst, reduce="sum")} es ``para cada nodo receptor,
suma los mensajes de sus aristas''. Es la versión rápida del sumatorio
$\bm m_j=\sum_{i\in\mathcal N(j)}\bm\epsilon_{ij}$. Con 8 bloques, la información
viaja hasta 8 aristas: ese es el ``radio de visión'' alrededor del arco.}

\subsection{El ensamblaje y el forward}

\texttt{MeshGraphNet.forward} encadena todo. La parte sutil es el \textbf{enrutado
del head}: si hay un head de conductancias (\texttt{full}/\texttt{enthalpy}), el
incremento se construye desde las \emph{latentes del processor} y \textbf{el decoder
se omite}; si no, se decodifica normalmente.

\srcref{src/models/meshgraphnet.py — MeshGraphNet.forward, l.708--732 (abreviado)}
\begin{lstlisting}
x_raw = x                                   # guardamos features crudas para el head
x, edge_attr = self.encoder(x, edge_attr)
for block in self.processor:                # L = 8 hops
    x, edge_attr = block(x, edge_index, edge_attr)

# heads de conductancias (full/enthalpy): el incremento sale de las latentes,
# NO del decoder (que se omite) -> garantiza la estructura GENERIC.
if self.generic_head is not None and self.generic_head.uses_node_latents:
    return self.generic_head(None, x_raw, batch=batch,
                             node_latents=x, edge_index=edge_index)

out = self.decoder(x)                        # baseline (sin GENERIC)
if self.generic_head is not None:            # head "energy-only" (proyeccion)
    out = self.generic_head(out, x_raw, batch)
return out
\end{lstlisting}

\begin{idea}[Por qué el decoder se ``puentea'']
En el GENERIC entálpico, el incremento \emph{no} puede ser un número libre por nodo
(eso violaría la degeneración $\Lw\bm 1=\bm0$). Tiene que salir de la
\textbf{Laplaciana de conductancias} actuando sobre $\mu=1/T$. Por eso, en ese modo,
el processor produce las \emph{latentes} con las que el head calcula las
conductancias, y el decoder libre se descarta (ni se instancia, l.669--670).
\end{idea}

%==============================================================================
\section{Normalización: por qué el head trabaja en kelvin físico}
\label{sec:norm}
%==============================================================================

La red se alimenta con features \textbf{normalizadas} (z-score) y emite un
$\Delta T$ \textbf{normalizado}. Las constantes se calculan una vez sobre todo el
corpus y se guardan en \texttt{stats.pt}:

\srcref{src/data/graph\_builder.py — NormalizeTransform, l.282--292}
\begin{lstlisting}
def __call__(self, data):
    x = data.x.clone(); m = self.mask                 # m = NORMALIZE_MASK
    x[:, m] = (x[:, m] - self.x_mean[m]) / self.x_std[m]   # z-score (one-hot intactos)
    data.x = x
    if data.y is not None:
        data.y = (data.y - self.y_mean) / self.y_std
    return data

def inverse_y(self, y):                               # vuelve a kelvin fisico
    return y * self.y_std + self.y_mean
\end{lstlisting}

Aquí aparece una tensión: el \textbf{head GENERIC necesita pensar en física}
(temperaturas en kelvin, $\mu=1/T$, energías), pero recibe features
\emph{normalizadas}. La solución: el head guarda sus propias copias de las
constantes (como \emph{buffers}, así viajan en el \texttt{state\_dict} y con
\texttt{.to(device)}) y \textbf{des-normaliza por dentro}. Se cargan una vez con
\texttt{set\_normalization}:

\srcref{src/models/meshgraphnet.py — \_GenericHeadBase, l.304--309 y \_phys\_col, l.338--343}
\begin{lstlisting}
self.register_buffer("x_mean", torch.zeros(f)); self.register_buffer("x_std", torch.ones(f))
self.register_buffer("y_mean", torch.zeros(1)); self.register_buffer("y_std", torch.ones(1))
# ...
def _phys_col(self, x_raw, idx):                 # des-normaliza la columna idx a kelvin
    c = x_raw[:, idx:idx+1]
    if bool(self.norm_mask[idx]):
        c = c * self.x_std[idx] + self.x_mean[idx]
    return c
\end{lstlisting}

\dummies{La red habla en ``unidades normalizadas''; la física habla en kelvin. El
head es el traductor: des-normaliza la temperatura para calcular $1/T$, hace toda la
termodinámica en kelvin, y al final vuelve a normalizar el $\Delta T$ para cumplir el
contrato de salida. Las constantes viajan dentro del modelo, así que un
\emph{checkpoint} reproduce la física exacta.}

En \texttt{train.py} esto se conecta con una sola línea tras crear el modelo:

\srcref{src/training/train.py — l.418--422}
\begin{lstlisting}
model = MeshGraphNet(cfg.model_config()).to(device)
if cfg.use_generic:
    model.set_normalization(stats)   # da al head las constantes z-score (en buffers)
\end{lstlisting}

%==============================================================================
\section{El GENERIC dentro de la red: conductancias y la tabla de entalpía}
\label{sec:generic}
%==============================================================================

Llegamos al núcleo: cómo \texttt{EnthalpyGenericThermalHead} implementa, línea a
línea, las ecuaciones del documento de formulación
\[
  \Delta H_i^{\mathrm{diss}}=\sum_{j\sim i} w_{ij}\Big(\tfrac1{T_i}-\tfrac1{T_j}\Big),
  \quad
  T_i^{\mathrm{new}}=H^{-1}\!\big(H(T_i)+\Delta H_i^{\mathrm{diss}}+\Delta H_i^{\mathrm{ext}}\big),
  \quad
  \Delta T_i=T_i^{\mathrm{new}}-T_i .
\]

\subsection{La jerarquía de \emph{heads}}

Hay tres heads, en herencia creciente (\texttt{generic\_mode}):
\texttt{"energy"} (solo degeneración, proyecta el incremento del decoder),
\texttt{"full"} (Laplaciana SPSD sobre \textbf{temperatura}) y \texttt{"enthalpy"}
(Laplaciana SPSD sobre \textbf{energía}; el correcto). \texttt{Enthalpy} hereda de
\texttt{Full} y solo cambia ``la variable que se conserva'' ($T\rightarrow h$) más
las tablas. El enrutado por modo está en el constructor del modelo:

\srcref{src/models/meshgraphnet.py — MeshGraphNet.\_\_init\_\_, l.669--679}
\begin{lstlisting}
conductance_head = cfg.use_generic and cfg.generic_mode in ("full", "enthalpy")
self.decoder = None if conductance_head else Decoder(cfg)   # se OMITE el decoder
if not cfg.use_generic:           self.generic_head = None
elif cfg.generic_mode == "enthalpy": self.generic_head = EnthalpyGenericThermalHead(cfg)
elif cfg.generic_mode == "full":     self.generic_head = FullGenericThermalHead(cfg)
else:                                self.generic_head = GenericThermalHead(cfg)
\end{lstlisting}

\subsection{Las conductancias aprendidas $w_{ij}\ge0$, simétricas}

La matriz $\bm M=\Lw$ se realiza como una Laplaciana de \textbf{conductancias
aprendidas}. Un MLP por arista produce un escalar no negativo (vía
\texttt{softplus}), multiplicado por una escala global aprendida
$e^{\alpha}$. El truco de la \textbf{simetría $w_{ij}=w_{ji}$ gratis}: el MLP solo
recibe features \emph{simétricas} bajo el intercambio $i\leftrightarrow j$
($\bm h_i+\bm h_j$, $(\bm h_i-\bm h_j)^2$, $T_i+T_j$, $|T_i-T_j|$).

\srcref{src/models/meshgraphnet.py — EnthalpyGenericThermalHead.forward, l.596--604}
\begin{lstlisting}
hi, hj = hlat[src], hlat[dst]
Tsrc, Tdst = T[src], T[dst]
feats = torch.cat([
    hi + hj,                       # simetrico bajo i<->j
    (hi - hj) ** 2,                # simetrico
    (Tsrc + Tdst) / (2.0 * tscale),
    (Tsrc - Tdst).abs() / tscale,
], dim=-1)
w = torch.exp(self.log_cond_scale) * F.softplus(self.cond_mlp(feats))   # (E,1) >= 0
\end{lstlisting}

\subsection{La acción de la Laplaciana, con \texttt{scatter}}

Con $\mu_i=1/T_i$, el incremento disipativo es la acción de la Laplaciana,
$(\Lw\mu)_i=\sum_{j\sim i}w_{ij}(\mu_i-\mu_j)$. En código: por cada arista dirigida
$s\!\rightarrow\!r$ se aporta $w\,(\mu_r-\mu_s)$ al receptor, y se suma con
\texttt{scatter}. Como las aristas son bidireccionales, esto es exactamente la
Laplaciana SPSD, y su suma por grafo es cero (energía conservada):

\srcref{src/models/meshgraphnet.py — EnthalpyGenericThermalHead.forward, l.606--612}
\begin{lstlisting}
mu = 1.0 / T                                            # gradiente de entropia dS/de = 1/T
# dH_diss = (L_w mu): Sigma_i dH_diss_i = 0  => 1er principio (energia conservada)
dH_diss = scatter(w * (mu[dst] - mu[src]), dst, dim=0,
                  dim_size=hlat.size(0), reduce="sum")
\end{lstlisting}

\begin{idea}[Ecuación $\leftrightarrow$ código, la traducción exacta]
\[
\underbrace{(\Lw\mu)_i=\!\!\sum_{j\sim i}\!w_{ij}(\mu_i-\mu_j)}_{\text{matemáticas}}
\;\;\equiv\;\;
\underbrace{\texttt{scatter(w*(mu[dst]-mu[src]), dst, reduce="sum")}}_{\text{PyTorch}}
\]
La conservación $\sum_i\Delta H_i^{\mathrm{diss}}=0$ \textbf{no se impone con una
penalización}: es una identidad algebraica de la Laplaciana (suma cero por columnas),
así que se cumple exactamente en cada \emph{forward}, gratis.
\end{idea}

\subsection{Las tablas que pasan de entalpía a temperatura}

Aquí está ``el rollo de las $c_p$'' hecho código. La curva $H(T)$ se \textbf{tabula
una sola vez} en el constructor, a partir del \emph{mismo} modelo de material que usa
el FEM (\texttt{MaterialProperties.cp\_apparent}), de modo que el mapa $T\leftrightarrow
H$ del sustituto es idéntico al de la verdad. Y se tabula en \textbf{unidades de
kelvin} ($H=h/(\rho\cpsens)$), no en J/m$^3$: ése es el reescalado de
acondicionamiento que evita que los gradientes se hundan bajo el $\varepsilon$ de Adam.

\srcref{src/models/meshgraphnet.py — EnthalpyGenericThermalHead.\_\_init\_\_, l.546--555}
\begin{lstlisting}
from simulation.thermal_solver import MaterialProperties
mat = MaterialProperties()
T_grid = np.linspace(self.T_FLOOR, self.T_TABLE_MAX, self.TABLE_POINTS)  # 200..8000 K, 4096 pts
cp_app = mat.cp_apparent(T_grid)                       # [J/(kg K)] con calor latente
dT = np.diff(T_grid, prepend=T_grid[0])
H_grid = np.cumsum((cp_app / mat.cp) * dT)             # H(T) en [K], MONOTONA
self.register_buffer("T_grid", torch.tensor(T_grid, dtype=torch.float32))
self.register_buffer("H_grid", torch.tensor(H_grid, dtype=torch.float32))
\end{lstlisting}

La \texttt{cp\_apparent} es las dos campanas gaussianas (fusión + vaporización) del
documento de formulación, ecuación del $\cpapp$:

\srcref{src/simulation/thermal\_solver.py — MaterialProperties.cp\_apparent, l.93--101}
\begin{lstlisting}
Tm = 0.5 * (self.T_solidus + self.T_liquidus); s = (self.T_liquidus - self.T_solidus)/4.0
dfl_dT = np.exp(-(((T - Tm)/s)**2)) / (s*np.sqrt(np.pi))                # campana fusion
s_v = self.vap_width/4.0
dfv_dT = np.exp(-(((T - self.T_vaporization)/s_v)**2)) / (s_v*np.sqrt(np.pi))  # campana vapor
return self.cp + self.latent_heat*dfl_dT + self.latent_heat_vap*dfv_dT  # = cp_app
\end{lstlisting}

¿Cómo se evalúa $H(T)$ y su inversa $T=H^{-1}(\cdot)$ \emph{dentro} del grafo de
autograd? Con una interpolación lineal a trozos \textbf{diferenciable}
(\texttt{\_interp1d}): busca el tramo con \texttt{searchsorted} y mezcla linealmente.
Su derivada es la pendiente local (finita), justo lo que necesita el backprop.

\srcref{src/models/meshgraphnet.py — \_interp1d, l.497--502}
\begin{lstlisting}
xq = x.clamp(xp[0], xp[-1]).squeeze(-1)
idx = torch.searchsorted(xp, xq.contiguous(), right=True).clamp(1, xp.numel()-1)
x0, x1 = xp[idx-1], xp[idx]; f0, f1 = fp[idx-1], fp[idx]
wgt = (xq - x0) / (x1 - x0).clamp_min(1e-12)
return (f0 + wgt*(f1 - f0)).unsqueeze(-1)
\end{lstlisting}

\dummies{Dos vectores, \texttt{T\_grid} y \texttt{H\_grid}, son la curva $h(T)$ del
otro documento (la rampa con dos rellanos). Ir de $T$ a $H$ es interpolar en una;
ir de $H$ a $T$ es interpolar en la otra. Como son \emph{buffers}, viajan en el
\emph{checkpoint}: el modelo lleva su propia física dentro.}

\subsection{El \emph{forward} completo del head entálpico}

Ahora se lee de un tirón, y es \textbf{exactamente} las tres ecuaciones del
encabezado de §\ref{sec:generic}: calcular la energía actual en $H$, sumar lo
disipativo y lo externo, y volver a $T$ invirtiendo la curva.

\srcref{src/models/meshgraphnet.py — EnthalpyGenericThermalHead.forward, l.633--640}
\begin{lstlisting}
# actualizacion de entalpia, y recuperar T por la curva de calor latente
H_cur = _interp1d(T, self.T_grid, self.H_grid)                 # H(T_i)
T_new = _interp1d(H_cur + dH_diss + dH_ext, self.H_grid, self.T_grid)  # H^{-1}(...)
dT_phys = T_new - T                                            # dT_i
return self._to_normalized_increment(dT_phys)                 # vuelve a unidades normalizadas
\end{lstlisting}

\begin{idea}[Por qué el calor latente sale ``gratis'']
En la meseta de fusión la \texttt{H\_grid} es casi vertical, así que la segunda
\texttt{\_interp1d} (de $H$ a $T$) devuelve un $T$ casi igual aunque le sumes mucho
$\Delta H$: un gran $\Delta H$ produce un $\Delta T$ minúsculo, \emph{igual que en el
FEM}. El ``freno $1/(\rho\cpapp)$'' del que hablaba el otro documento es,
literalmente, la pendiente de la curva tabulada.
\end{idea}

\subsection{La fuente enriquecida}

La parte disipativa solo enfría; quien calienta el pico es la \textbf{fuente
externa} $\Delta H^{\mathrm{ext}}$. En la versión enriquecida, las ganancias
escalares se modulan por nodo con dos MLPs, \textbf{inicializados a cero} (así
arranca idéntico a la versión escalar), con \texttt{softplus}$\ge0$ y
\emph{gateado} por la fuente Goldak $q$ (no puede inventar energía lejos del arco):

\srcref{src/models/meshgraphnet.py — \_\_init\_\_ (zero-init), l.569--575 y forward, l.625--629}
\begin{lstlisting}
self.source_mlp = build_mlp(hdim+2, hdim, 1, ...); self.cool_mlp = build_mlp(hdim+2, hdim, 1, ...)
for m in (self.source_mlp, self.cool_mlp):
    nn.init.zeros_(m[-1].weight); nn.init.zeros_(m[-1].bias)   # arranca = ganancia escalar
# ... en el forward:
src_gain  = F.softplus(self.source_gain + self.source_mlp(torch.cat([hlat, q_n, T_n], -1)))
cool_gain = F.softplus(self.cool_gain   + self.cool_mlp(torch.cat([hlat, T_n, tinf_n], -1)))
dH_ext = src_gain * q + cool_gain * robin * (t_inf - T)        # >=0 calienta / signo-correcto enfria
\end{lstlisting}

Nótese que \texttt{dH\_diss} \textbf{no se toca}: sigue siendo estructura-preservante.
Solo el término externo (que de todos modos no está sujeto a las garantías) gana
flexibilidad. Energía interna conservada y $\dd S/\dd t\ge0$ se mantienen.

%==============================================================================
\section{El entrenamiento: datos, ventanas y push-forward}
\label{sec:train}
%==============================================================================

\subsection{El \texttt{Dataset} y el \emph{cache} en disco}

\texttt{WeldingGraphDataset} es un \texttt{Dataset} de PyG perezoso: procesa los
\texttt{.npz} crudos a un \texttt{.pt} por par de fotogramas, calcula
\texttt{stats.pt} (las constantes de normalización, en \emph{streaming} para no
cargar todo en RAM) y mantiene un \textbf{índice determinista}
\texttt{\_index = [(fname, sim\_idx, t)]} ordenado por $(\text{sim},t)$. Ese índice
es la clave que luego usan el \emph{split} y las ventanas.

\srcref{src/data/graph\_builder.py — process() (acumuladores) l.383--394 y get() l.421--426}
\begin{lstlisting}
xn = data.x.numpy().astype(np.float64)
x_sum += xn.sum(0); x_sqsum += (xn**2).sum(0); x_count += xn.shape[0]   # media/var en streaming
torch.save(data, processed_dir / f"data_{sim_idx}_{t}.pt")             # un .pt por grafo
# ...
def get(self, idx):
    fname, sim_idx, t = self._index[idx]
    return torch.load(processed_dir / f"data_{sim_idx}_{t}.pt", weights_only=False)
\end{lstlisting}

\subsection{Split por simulación (no por grafo)}

Crucial en SciML: si parties por grafo, fotogramas casi idénticos del mismo cordón
caen en train y val, e inflas el resultado. Aquí se baraja la \textbf{lista de
simulaciones} (con semilla) y \textbf{cada grafo hereda el fold de su simulación}.

\srcref{src/training/utils.py — split\_by\_simulation, l.180--188 (abreviado)}
\begin{lstlisting}
for graph_idx, sim in enumerate(ids):          # ids[graph] = simulacion de ese grafo
    if   sim in train_sims: train_idx.append(graph_idx)
    elif sim in val_sims:   val_idx.append(graph_idx)
    else:                   test_idx.append(graph_idx)
\end{lstlisting}

\subsection{Las ventanas del push-forward}

Para entrenar con $K$ pasos hace falta servir \textbf{ventanas de $K$ grafos
consecutivos de la misma simulación}. \texttt{WindowedSubset} precalcula los
arranques válidos (mismo \texttt{sim\_idx} en los dos extremos $\Rightarrow$ los $K$
del medio son consecutivos):

\srcref{src/training/utils.py — WindowedSubset, l.425--439 (abreviado)}
\begin{lstlisting}
for i in range(n - self.k + 1):
    if int(index[i][1]) not in train: continue          # arranque en sim de train
    if int(index[i + self.k - 1][1]) == int(index[i][1]):   # K consecutivos sin cruzar sim
        self.starts.append(i)
# ...
def __getitem__(self, j):
    i = self.starts[j]
    return [self.base.get(i + k) for k in range(self.k)]   # lista de K grafos crudos
\end{lstlisting}

Y \texttt{collate\_windows} reorganiza un \emph{batch} de $B$ ventanas en $K$
\textbf{mini-batches alineados por paso}: el batch $k$ junta el $k$-ésimo grafo de
cada ventana. Como dentro de una ventana el orden de nodos es idéntico (misma
simulación) y las ventanas se concatenan en el mismo orden en cada $k$, \textbf{el
nodo $j$ del batch $k$ es el nodo $j$ del batch $k{+}1$}: por eso un único vector
\texttt{T\_hat} se puede ``rodar'' por toda la ventana.

\srcref{src/training/utils.py — collate\_windows, l.453--454}
\begin{lstlisting}
k = len(batch[0])
return [Batch.from_data_list([window[step] for window in batch]) for step in range(k)]
\end{lstlisting}

\dummies{Una ``ventana'' son $K$ fotogramas seguidos del mismo cordón. El
\emph{collate} los apila de forma que el nodo $j$ siempre es el mismo nodo físico a
lo largo de los $K$ pasos. Eso permite realimentar la predicción paso a paso dentro
del batch, que es justo lo que pide el push-forward.}

\subsection{El bucle de una época: push-forward y la pérdida con $K$}

Esta es la función estrella, y es el \textbf{Algoritmo de push-forward} del paper.
Se siembra $\hat T$ con el campo verdadero (más ruido) al inicio de la ventana; se
rueda el modelo $K$ pasos \textbf{realimentando su propia predicción} (sin
\texttt{detach}, para que el gradiente atraviese todo el desenrollado: BPTT); y se
acumula el MSE en temperatura física escalado por $\sigma_T$.

\srcref{src/training/train.py — \_train\_one\_epoch, l.286--309 (abreviado)}
\begin{lstlisting}
# 1) sembrar T desde el inicio de la ventana, con ruido proporcional a la temperatura
T_hat = dynamic_temperature_noise(b0.x[:, TEMPERATURE_INDEX].clone(),
                                  b0.x[:, T_INF_INDEX], beta=beta, floor=floor)
loss = b0.x.new_zeros(())
for bk in steps:                                  # K pasos
    x = bk.x.clone(); x[:, TEMPERATURE_INDEX] = T_hat   # realimenta nuestra prediccion
    x_norm = x.clone()
    x_norm[:, mask] = (x[:, mask] - nt.x_mean[mask]) / nt.x_std[mask]   # normaliza
    dT = nt.inverse_y(model(x_norm, bk.edge_index, bk.edge_attr, batch=bk.batch)).squeeze(-1)
    T_pred = T_hat + dT
    T_true = bk.x[:, TEMPERATURE_INDEX] + bk.y[:, 0]          # T verdadera del paso k
    loss = loss + (((T_pred - T_true) / t_scale) ** 2).mean() # MSE escalado por sigma_T
    T_hat = T_pred                                            # feed-forward (SIN detach -> BPTT)
loss = loss / len(steps)
loss.backward()
if grad_clip > 0: clip_grad_norm_(model.parameters(), grad_clip)
optimizer.step()
\end{lstlisting}

Esto implementa
$\mathcal L=\frac1K\sum_{k=0}^{K-1}\big\lVert(T^{(k)}_{\mathrm{pred}}-T^{(k)}_{\mathrm{true}})/\sigma_T\big\rVert^2$.
Con \texttt{pushforward\_steps=1} se recupera el entrenamiento clásico de un paso;
el headline usa $K=2$. El ruido lo pone \texttt{dynamic\_temperature\_noise} (mayor
cerca del baño, $\sim0$ en el campo lejano):

\srcref{src/training/utils.py — dynamic\_temperature\_noise, l.385--389}
\begin{lstlisting}
sigma = beta * (T - T_inf).abs() + floor       # grande cerca de la fusion, ~0 lejos
eta = torch.randn(T.shape, ...) * sigma
return T + eta
\end{lstlisting}

\begin{idea}[La línea más importante de todo el entrenamiento]
\texttt{T\_hat = T\_pred} \emph{sin} \texttt{detach}. Si pusieras \texttt{.detach()}
ahí, cada paso entrenaría aislado (teacher forcing) y volverías al objetivo de un
paso, que \emph{no} está alineado con el rollout largo. No hacer detach es lo que
propaga el gradiente a través de los $K$ pasos (BPTT) y alinea el entrenamiento con
la inferencia. Esa única ausencia de \texttt{detach} \emph{es} el push-forward.
\end{idea}

\subsection{Validación = rollout autoregresivo (no un paso)}

La métrica que selecciona el \emph{checkpoint} \textbf{no} es la pérdida de un paso:
es el RMSE de un \textbf{rollout autoregresivo completo} sobre las simulaciones de
validación (el modelo despliega toda la trayectoria realimentándose). Por eso vive
en \texttt{rollout.py} y se invoca así:

\srcref{src/training/train.py — \_validate\_rollout, l.333--340}
\begin{lstlisting}
for name, result in val_results.items():
    rollout = run_autoregressive_rollout(model, result, normalizer, device=device)
    per_sim[name] = rollout.rmse
per_sim["mean"] = float(np.mean(list(per_sim.values())))
\end{lstlisting}

\subsection{El bucle exterior: optimizador, \emph{scheduler}, \emph{checkpoints}}

El resto de \texttt{train()} es PyTorch ``de toda la vida'', con detalles pensados
para \emph{pods} que se pueden cortar: \texttt{AdamW} (lr $10^{-3}$, weight decay
$10^{-4}$), \emph{scheduler} configurable (\texttt{plateau} en el headline),
guardado \textbf{atómico} de \texttt{best\_model.pt} (mejor métrica) y
\texttt{last\_model.pt} (estado completo para \texttt{--resume}). El guardado atómico
(temp $+$ \texttt{os.replace}) evita dejar un checkpoint a medias si el pod muere.

\srcref{src/training/train.py — selección de mejor y guardado, l.490--505 (abreviado)}
\begin{lstlisting}
improved = metric < best_metric - cfg.early_stop_min_delta
if improved:
    best_metric = metric; best_epoch = epoch; epochs_since_improve = 0
    _atomic_save({"model_state": model.state_dict(), "config": asdict(cfg),
                  "epoch": epoch, "metric": best_metric, ...}, ckpt_dir / "best_model.pt")
\end{lstlisting}

%==============================================================================
\section{Inferencia: el rollout autoregresivo en detalle}
\label{sec:rollout}
%==============================================================================

En inferencia el modelo recibe \textbf{solo} el campo inicial $T^0$ y la trayectoria
Goldak (conocida). Como la trayectoria es determinista, \textbf{todas las features
dependientes del tiempo se precalculan} para todos los pasos con el mismo
\texttt{build\_graph\_sequence} del entrenamiento (paridad total). En el bucle se
\textbf{sobrescribe solo la columna de temperatura} con la predicción que se va
acumulando:

\srcref{src/training/rollout.py — run\_autoregressive\_rollout, l.188--212 (abreviado)}
\begin{lstlisting}
graphs = build_graph_sequence(result, sim_id=sim_id)     # MISMO builder que en train
feats = torch.stack([g.x for g in graphs]).to(device)    # (S-1, N, 12) precomputado
edge_index = graphs[0].edge_index; edge_attr = graphs[0].edge_attr   # topologia fija

T_cur = gt[0].clone()                                    # arranca de la T^0 verdadera
for t in range(n_steps):
    x = feats[t].clone(); x[:, TEMPERATURE_INDEX] = T_cur    # sobrescribe SOLO la T
    x_norm = x.clone(); x_norm[:, mask] = (x[:, mask] - nt.x_mean[mask]) / nt.x_std[mask]
    dT = nt.inverse_y(model(x_norm, edge_index, edge_attr)).squeeze(-1)
    T_cur = T_cur + dT                                       # autoregresivo
    history.append(T_cur.clone())
\end{lstlisting}

\dummies{Inferencia $=$ ``dame $T^0$ y te doy toda la película''. La red predice
$\Delta T$, lo suma, y vuelve a meter esa $T$ como entrada del paso siguiente.
Solo la temperatura se realimenta; el arco y la geometría ya se conocen. Es
\emph{idéntico} al bucle interno del push-forward, pero sin gradiente y para los
$\sim$1500 pasos completos.}

\begin{idea}[El mismo bucle, dos usos]
El \emph{feed-forward} ``sobrescribe T $\rightarrow$ normaliza $\rightarrow$ modelo
$\rightarrow$ \texttt{inverse\_y} $\rightarrow$ $T{+}{=}\Delta T$'' aparece
\textbf{dos veces}, idéntico: en \texttt{rollout.py} (inferencia, \texttt{no\_grad},
$\sim$1500 pasos) y dentro de \texttt{\_train\_one\_epoch} (entrenamiento,
diferenciable, $K$ pasos). Esa simetría es la que alinea objetivo y despliegue.
\end{idea}

%==============================================================================
\section{Cómo responder (casi) cualquier pregunta del repo}
\label{sec:faq}
%==============================================================================

Con el mapa de §\ref{sec:mapa} y las traducciones ecuación$\leftrightarrow$código,
estas preguntas típicas se responden ``de memoria'':

\begin{description}[leftmargin=0em,style=nextline]
\item[``¿Dónde se garantiza la conservación de energía?'']
En la suma cero por columnas de la Laplaciana:
\texttt{scatter(w*(mu[dst]-mu[src]), dst, "sum")} (meshgraphnet.py l.611). No hay
penalización; es identidad algebraica.
\item[``¿Dónde está el calor latente?'']
En \texttt{MaterialProperties.cp\_apparent} (thermal\_solver.py l.93), tabulado a
\texttt{H\_grid} en el \texttt{\_\_init\_\_} del head (meshgraphnet.py l.553), y
aplicado al invertir con \texttt{\_interp1d} (l.635).
\item[``¿Por qué la red no usa coordenadas?'']
Porque \texttt{NODE\_FEATURE\_NAMES} (graph\_builder.py l.48) solo tiene features
relativas; \texttt{pos} se guarda aparte (l.257) y nunca entra en \texttt{x}.
\item[``¿Qué hace exactamente $K=2$?'']
Desenrolla 2 pasos realimentando la predicción \emph{sin} \texttt{detach}
(train.py l.303) y promedia el MSE escalado (l.302,304).
\item[``¿Por qué el head trabaja en kelvin si la red está normalizada?'']
Porque guarda las constantes en \emph{buffers} y des-normaliza con \texttt{\_phys\_col}
(meshgraphnet.py l.338); se cargan con \texttt{set\_normalization} (train.py l.421).
\item[``¿Cómo se elige el mejor modelo?'']
Por el RMSE de \textbf{rollout} de validación, no por la loss de un paso
(train.py l.466--505); \texttt{best\_model.pt} guarda ese mínimo.
\item[``¿Por qué el decoder a veces no existe?'']
En modo \texttt{full}/\texttt{enthalpy} el incremento sale de las latentes
(estructura GENERIC), así que el decoder se pone a \texttt{None} (meshgraphnet.py
l.669--670) y el forward enruta por el head (l.720).
\end{description}

%==============================================================================
\section{Recapitulación: ecuación $\rightarrow$ archivo $\rightarrow$ línea}
\label{sec:resumen}
%==============================================================================

\renewcommand{\arraystretch}{1.25}
\noindent\small
\begin{tabular}{@{}p{0.36\linewidth}p{0.60\linewidth}@{}}
\toprule
Concepto / ecuación & Implementación (archivo · líneas) \\
\midrule
12 features de nodo, target $\Delta T$ & \texttt{graph\_builder.py} · l.48--73, 246--260 \\
Encoder $\bm v^{(0)},\bm\epsilon^{(0)}$ & \texttt{meshgraphnet.py} · Encoder l.170--173 \\
Paso de mensajes (update + scatter + residual) & \texttt{meshgraphnet.py} · GraphNetBlock l.209--224 \\
Enrutado del head / decoder puenteado & \texttt{meshgraphnet.py} · l.669--670, 708--732 \\
Normalización z-score, \texttt{inverse\_y} & \texttt{graph\_builder.py} · l.282--292 \\
Head en kelvin (buffers, \texttt{\_phys\_col}) & \texttt{meshgraphnet.py} · l.304--343 \\
Conductancias $w_{ij}\ge0$ simétricas & \texttt{meshgraphnet.py} · l.596--604 \\
$(\Lw\mu)_i$ vía \texttt{scatter} (1er principio) & \texttt{meshgraphnet.py} · l.606--612 \\
$\cpapp$ (campanas de calor latente) & \texttt{thermal\_solver.py} · l.93--101 \\
Tabla $H(T)$ en kelvin (reescalado) & \texttt{meshgraphnet.py} · l.546--555 \\
Interp. diferenciable $H\leftrightarrow T$ & \texttt{meshgraphnet.py} · \_interp1d l.497--502 \\
Update entálpico $T{=}H^{-1}(H{+}\Delta H)$ & \texttt{meshgraphnet.py} · l.633--640 \\
Fuente enriquecida (zero-init, softplus, gate) & \texttt{meshgraphnet.py} · l.569--575, 625--629 \\
Split por simulación & \texttt{utils.py} · l.180--188 \\
Ventanas $K$ + collate alineado & \texttt{utils.py} · l.425--454 \\
Ruido dinámico & \texttt{utils.py} · l.385--389 \\
Push-forward $K$ + pérdida (BPTT) & \texttt{train.py} · l.286--309 \\
Validación por rollout & \texttt{train.py} · l.333--340 ; \texttt{rollout.py} · l.188--212 \\
Best/last checkpoint atómico, resume & \texttt{train.py} · l.324--329, 441--523 \\
\bottomrule
\end{tabular}

\bigskip
\begin{idea}[La moraleja de implementación]
La formulación se traduce en \textbf{tres patrones de código} que reconocerás en
todas partes: (1) \texttt{scatter} sobre aristas $=$ ``operador de grafo'' (paso de
mensajes \emph{y} Laplaciana GENERIC); (2) dos \texttt{\_interp1d} contra las tablas
\texttt{T\_grid}/\texttt{H\_grid} $=$ ``el calor latente y la conversión
$T\leftrightarrow H$''; (3) el bucle ``sobrescribe T $\rightarrow$ modelo
$\rightarrow$ $T{+}{=}\Delta T$'', con gradiente ($K$ pasos) en entrenamiento y sin
él ($\sim$1500) en inferencia. Si entiendes esos tres patrones, entiendes el repo.
\end{idea}

\end{document}
